\begin{figure}[htbp]
  \centering
  \begin{tikzpicture}[x=10mm,y=10mm,alt={A full seven-byte ring has head and tail both five and used seven. Physical indices zero through six contain decimal bytes 66, 67, 2, 120, 121, 3, 65. The oldest record begins with length three at index five, then payload A at six, B at zero and C at one. The second record has length two at index two and payload x and y at three and four. Shading marks the first record; a wrap arrow connects the physical ends.}]
    \foreach \i in {0,1,5,6} {\fill[black!18] (\i,0) rectangle ++(1,1);}
    \foreach \i in {0,...,7} {\draw (\i,0) -- (\i,1);}
    \draw (0,0) -- (7,0) (0,1) -- (7,1);
    \foreach \i/\value/\kind in {0/66/B,1/67/C,2/2/length,3/120/x,4/121/y,5/3/length,6/65/A} {
        \node[font=\small,anchor=south] at (\i+0.5,1.05) {\i};
        \node[font=\small] at (\i+0.5,0.65) {\value};
        \node[font=\footnotesize] at (\i+0.5,0.25) {\kind};
      }
    \draw[-{Stealth},rounded corners=2pt] (7,0.5) -- (7.4,0.5) -- (7.4,1.7) -- (-0.4,1.7) -- (-0.4,0.5) -- (0,0.5);
    \draw[-{Stealth}] (5.5,-0.7) -- (5.5,0);
    \node[anchor=north,font=\small] at (5.5,-0.75) {$head=tail=5$};
    \node[anchor=north,font=\small] at (2,-0.2) {$used=7$ (full)};
  \end{tikzpicture}
  \caption{Logical record order crosses the physical boundary. Each length prefix occupies a byte of capacity, and equal head and tail indices here describe a full queue.}
  \label{fig:revision-ring-full}
\end{figure}
